\hypertarget{discussion}{%
\chapter{Discussion}\label{discussion}}

This chapter interprets the empirical findings of Chapter 5 in the context of the broader research questions identified in Chapter 1 and the literature reviewed in Chapter 3. The discussion is organised around six interpretive lenses: the consolidation of the key quantitative findings into a coherent claim about the system's behaviour (6.1), the architectural strengths that the evaluation supports (6.2), the challenges encountered and the specific design decisions taken to address them (6.3), the implications of the results for food safety practice and regulatory compliance (6.4), the comparative position of the IDRISK2 system relative to the state of the art (6.5), and an honest accounting of the limitations of the current evaluation (6.6).

\hypertarget{interpretation-of-key-findings}{%
\section{Interpretation of Key Findings}\label{interpretation-of-key-findings}}

The single most consequential finding of the evaluation is that the IDRISK2 system's classification accuracy is entirely attributable to the Retrieval-Augmented Generation layer. The Vision-Language Model, evaluated in isolation, produced zero correct base-term assignments across 34 images; the full RAG-augmented pipeline produced 10 correct assignments and 20 family-level assignments. The marginal contribution of the retrieval and reranking stages is therefore not an incremental improvement over a baseline but the entire useful signal of the system. This result transforms the methodological status of the RAG architecture from an enhancement to a necessary precondition for the system's existence as a regulatory-grade classifier.

The pattern of VLM failures provides a fine-grained characterisation of the hallucination phenomenon in the structured-taxonomy domain. The model did not produce malformed outputs: every preliminary code conformed to the FoodEx2 syntactic format of an A-prefix followed by four alphanumeric characters, and every accompanying natural-language description of the image was substantively accurate. The failure mode is specifically the binding between visual content and taxonomy identifier, where the model lacks reliable parametric knowledge of the 31,303-term taxonomy and resorts to plausible-looking fabrications. The recurrence of a small set of codes as fallbacks (A0FQJ five times, A0EKC four times) is consistent with the model having memorised a handful of FoodEx2-shaped strings during training and reusing them as generic answers when no confident parametric mapping exists. Crucially, these fabrications are emitted with high confidence in 33 of 34 cases, which renders the confidence signal of the upstream VLM useless for any downstream calibration purpose.

The cross-VLM evaluation reported in Section 5.7.2 strengthens this characterisation from a single-model observation to a class-level claim, and further refines it through comparison with the open-source 7B-parameter alternative. When the same 34-image dataset is processed with Claude Sonnet 4.6 instead of GPT-4o, the bare-VLM exact-match accuracy remains at 0.0\% and the end-to-end system accuracy remains at 29.4\%; the two proprietary VLMs converge on identical headline metrics. When the substitution extends to LLaVA-NeXT-Mistral-7B, the bare-VLM exact-match accuracy remains at 0.0\% but the end-to-end system accuracy collapses to 2.9\%. The triple convergence of the bare-VLM result at zero across architectures of widely differing parameter count and training regime makes the hallucination phenomenon a property of the contemporary Vision-Language Model class as a whole, not a quirk of any single model nor of the proprietary subclass; it elevates the methodological status of the retrieval-grounded correction from a useful mitigation to an empirically necessary one across the available alternatives. The same comparison delimits the architectural substitutability of the VLM layer that Section 6.2 had previously characterised as a design-time property: substitutability holds robustly within the proprietary class, where GPT-4o and Claude yield equivalent system accuracy, but breaks down at the proprietary-to-7B-open-source boundary, where system accuracy drops by an order of magnitude. The determining factor in system accuracy is therefore the joint quality of the VLM's preliminary signal and the retrieval and reranking stages downstream of it, not either component in isolation.

Conversely, the pattern of RAG successes confirms the design choice documented in Section 4.5.1 to index the FoodEx2 taxonomy at term granularity rather than as further-chunked passages. In every classification, all five retrieved candidates were atomic taxonomy term nodes, and the LLM-based generator selected from this candidate set in the 30 cases that did not exhibit label leak. The brigadeiro case study in Section 5.8 is the cleanest instance of this dynamic: the VLM proposed a hallucinated code in the unrelated taxonomic family of "Lime infusion flowers", but the retriever surfaced "Chocolate coated confectionery" as the top reranked candidate, and the generator correctly overrode the VLM's proposal. The RAG layer, in this functioning regime, operates exactly as the literature on hallucination mitigation predicts: by replacing parametric guessing with verifiable evidence retrieval.

The family-level accuracy of 58.8\% relative to the exact-match accuracy of 29.4\% indicates that the dominant residual failure mode of the system, after the RAG correction takes effect, is granularity selection within the correct hierarchy branch rather than gross taxonomic misclassification. Of the 24 cases where the system did not exactly match the ground truth, 10 fell within the same three-character taxonomy prefix as the ground truth and 14 did not. The granularity errors are typically distinctions that a non-expert human annotator would also find difficult, for example A00BG "Chocolate cake" versus A00BF "Chocolate-based cakes" for the brownie image, or A025N "Sliceable or firm cooked sausages" versus A025C "Chorizo and similar" for the chouriço image. These are tractable improvements through better prompting and richer term descriptions in the indexed corpus. The remaining 14 cases of cross-family error are the substantively harder problem and motivate the architectural changes proposed in Section 8.1.

\hypertarget{advantages-of-the-proposed-idrisk2-system}{%
\section{Advantages of the Proposed IDRISK2 System}\label{advantages-of-the-proposed-idrisk2-system}}

The evaluation supports three architectural choices made in the system's design. The first is the decision to combine the EFSA technical PDFs with the Appendix B taxonomy XLSX as two complementary indexing streams with heterogeneous chunking strategies. The PDFs supplied governance and methodology context useful for grounding generation, while the atomic taxonomy term nodes supplied the candidate identifiers that the classifier required to produce correct outputs. The cross-encoder consistently reranked term nodes above PDF chunks across all 34 evaluations, which validates the heterogeneous-chunking design: the choice to leave each taxonomy term as an unfragmented atomic node preserves the semantic locality that the reranker depends on, and the choice to chunk PDFs semantically preserves the descriptor definitions intact for the generation context.

The second supported choice is the use of an Advanced RAG architecture combining pre-retrieval query rewriting, HyDE-based dense retrieval, sparse BM25 lexical retrieval, Reciprocal Rank Fusion, and ColBERT cross-encoder reranking. The mean top-1 rerank score of 0.693 is substantially above the 0.611 ceiling observed in earlier pilot runs that indexed only the EFSA PDFs, confirming the value of the layered retrieval pipeline. The fact that 100\% of retrieved candidates were taxonomy terms across all evaluations indicates that the cross-encoder discriminates effectively between meta-level documentation and dish-level taxonomy entries, which is the discrimination most directly relevant to classification.

The third supported choice is the closed, self-hosted deployment architecture documented in Section 4.6. The BAAI/bge-m3 embedding model, the Chroma vector store, the cross-encoder reranker, and the rule-based security infrastructure all operate locally on infrastructure controlled by the deploying organisation. Only the calls to the proprietary VLM (GPT-4o or Claude, depending on configuration) exit the deployment boundary; the architecture supports substitution of either proprietary VLM by an open-source LLaVA-NeXT deployment at the cost of additional local compute and the loss of state-of-the-art reasoning quality. This design preserves the option for fully self-hosted operation as the open-source VLM ecosystem matures, without requiring re-architecting of the retrieval or generation layers.

Beyond accuracy considerations, the evaluation supports the operational feasibility of the system. The median latency of 22.3 seconds with a standard deviation of 1.5 seconds across heterogeneous inputs indicates predictable per-classification cost, which is a precondition for inclusion in time-sensitive inspection workflows. The audit trail produced by every classification, comprising the cryptographically chained HMAC entries written by the AuditLogger and the per-decision provenance records written by the generation component, provides the documentation required by the EU AI Act for high-risk AI systems, regardless of the classification's correctness. The system therefore meets the documentation and traceability obligations of the regulatory framework even where it fails to meet the accuracy obligations, which is an important separation for deployment in supervised regulatory contexts.

\hypertarget{challenges-encountered-and-solutions-adopted}{%
\section{Challenges Encountered and Solutions Adopted}\label{challenges-encountered-and-solutions-adopted}}

Four substantive architectural challenges emerged during the development and evaluation of the system. Each is documented here with the design response adopted and the residual risk that remains for future iterations.

The first challenge is the label leak phenomenon observed in 4 of 34 cases under the GPT-4o configuration (11.8\% of the evaluation set), 8 of 34 cases under the Claude configuration (23.5\%), and 3 of 34 cases under the LLaVA-NeXT-7B configuration (8.8\%). When the VLM emits a preliminary FoodEx2 code together with a fabricated natural-language label, the final structured generator can carry both the code and the label through to its output even when the retrieved context does not support the assignment. This failure mode arises because the generation prompt presents both the VLM's preliminary classification and the retrieved candidates as joint context, and the LLM occasionally treats the VLM's claim as a permitted continuation rather than a hypothesis to be checked against retrieval. The mitigating design response is the explicit instruction in the generation system prompt to base classification decisions "only on the retrieved documentation provided" (Section 4.5.3). The interpretation of the leak rate must be read jointly with downstream accuracy. Under the open-source LLaVA configuration the leak rate is the lowest of the three VLMs, but this apparent advantage is the symptomatic outcome of LLaVA's near-uniform emission of a single generic-fallback code (A0FQJ "Mixed dish with meat and vegetables" in 33 of 34 cases), which the retriever and the cross-encoder readily override, rather than evidence of strong VLM-classifier discipline. The leak metric therefore diagnoses VLM-classifier interaction and is not a standalone quality indicator. The residual risk is documented as an architectural defect requiring a downstream validation step in which the emitted code must be verified against the retrieved candidate set before the output is returned; the cross-VLM evidence strengthens the case for this mitigation by demonstrating that VLM substitution alone does not eliminate the failure mode, that under some VLM choices it amplifies it (Claude), and that under others it merely shifts the dominant failure mode from leak to retrieval contamination (LLaVA).

The second challenge is the cultural-coverage limitation of the FoodEx2 taxonomy itself. The dataset's emphasis on Portuguese and Brazilian dishes surfaced repeated cases in which no specific term exists for the food being classified: acarajé (no specific code, mapped to "Legumes based dishes"), brigadeiro (no specific code, mapped to "Chocolate coated confectionery"), pastel (no specific code, mapped to "Sandwiches, pizza and other stuffed bread-like cereal products"), rabanada (no specific code, mapped to the generic base term "Wheat bread and rolls, white"), and the Brazilian cuscuz nordestino (no specific code, mapped to the structurally similar "Couscous" entry derived from Maghrebi cuisine). This is a coverage limitation of the EFSA taxonomy rather than a defect of the IDRISK2 system, but it materially affects the achievable accuracy ceiling on the Lusophone evaluation set. The mitigating response, adopted at ground-truth annotation time, was to apply the FoodEx2-prescribed practice of selecting the closest generic base term and intending that the dish-specific characterisation be carried by facets. The residual risk is that the system's outputs in these cases convey less information than a Brazilian or Portuguese inspector would expect, which raises an open question about whether the European-centric taxonomy is appropriate as the unique target ontology for inspection workflows in Brazil and Portugal.

The third challenge is the cascading failure pattern observed when the upstream VLM misidentifies the visual content of the image. The acarajé case discussed in Section 5.8 is the prototypical example: a single error in the VLM's visual description ("fried plantain" instead of "fried bean fritter") propagates through query rewriting, HyDE generation, retrieval, and reranking, with each downstream stage operating faithfully on the contaminated input. The system's modular design treats each stage as a trusted producer for the next, and there is no mechanism to detect or rollback an upstream error. The mitigating response adopted in the current architecture is the chain-of-thought reasoning field in the structured generation output, which produces a human-readable record of the system's classification rationale that an auditor can inspect after the fact. The residual risk is that errors of this category are not detected automatically and that their detection depends on human review, which the evaluation showed to be ineffective at 4.2\% recall on actual errors.

The fourth challenge is the calibration of the system's confidence signal. The evaluation revealed a false high-confidence rate of 71.9\% and a silent error rate of 95.8\%, figures that establish the current confidence mechanism as effectively non-functional for regulatory purposes. The mitigating design response was the introduction of the configurable REVIEW\_TRIGGER\_LEVELS policy (Section 4.5.3) that flags any classification with low or medium confidence for human review by default. This policy correctly triggered review on 2 of the 34 evaluations, but in only one case (mandioca) was the flagged classification incorrect, while 23 incorrect classifications passed through with high confidence and no flag. The residual risk is substantial: under the current policy, the human-review mechanism does not satisfy the human-oversight obligation of EU AI Act Article 14 for high-risk AI applications, and the architectural redesign of the trigger is documented as priority future work in Section 8.1.

\hypertarget{implications-for-food-safety-and-compliance}{%
\section{Implications for Food Safety and Compliance}\label{implications-for-food-safety-and-compliance}}

The deployment of the IDRISK2 system as an autonomous classifier is not supported by the current evaluation. An exact-match accuracy of 29.4\% combined with a 71.9\% false high-confidence rate places the system far below the threshold at which automatic classifications could replace human judgement in a regulatory food-safety workflow. The implication for deployment is therefore that the system, in its present configuration, must operate as a decision-support tool with mandatory human review on every classification, not only on those flagged by the automatic trigger. This reverses the implicit assumption of the architecture, in which review is the exception rather than the rule, and brings the operational practice into line with the calibration that the evaluation has actually measured.

Under that revised deployment model, the system's contribution to inspection workflows is best framed as the production of a high-quality candidate classification and a complete audit trail, both of which are presented to a human inspector for final adjudication. In this framing, the family-level accuracy of 58.8\% is more directly relevant than the exact-match accuracy: 58.8\% of the time, the system places the inspector within the correct three-character branch of the taxonomy and reduces the search space for human review from 31,303 terms to a small set of related candidates. The inspector retains responsibility for the final code selection, the system's confidence signal becomes an advisory rather than authoritative input, and the audit trail produced by every classification supports post-hoc review by supervisors or auditors. This framing is consistent with the human-in-the-loop deployment model recommended for high-risk AI systems by the EU AI Act and is the model the IDRISK2 system can support today. The web frontend described in Section 4.8 was designed against this deployment model: the role-specific screens surface the candidate classification, the live pipeline progress, and the audit trail to the inspector and supervisor at the appropriate points in the workflow, while the human-review queue on the supervisor dashboard materialises the human-oversight mechanism on which the EU AI Act compliance posture depends.

With respect to the GDPR, the system's compliance posture is robust. The closed deployment architecture documented in Section 4.6 keeps all product data on infrastructure controlled by the deploying authority, the data-minimisation mechanism discards raw image bytes immediately after feature extraction and persists only the SHA-256 hash and the classification result, and the chained HMAC audit log provides the tamper-evident traceability required by Article 22 for automated decision-making. The label leak phenomenon discussed in Section 6.3 is, in this context, a defect specifically of explainability: when the system emits a label that does not correspond to the official taxonomy entry for the assigned code, the audit record contains a misleading rationale that cannot be unambiguously reconciled with the FoodEx2 ground truth. This is a tractable defect, adding a final validation step that resolves every emitted code against the taxonomy and rejects mismatches, but its mitigation is required before the system can claim full compliance with the explainability obligations of Article 22.

With respect to the EU AI Act, the system's current calibration failure is the dominant compliance gap. Article 14's human-oversight requirement presupposes a meaningful trigger that escalates uncertain decisions to human review; a trigger that captures 4.2\% of actual errors does not meet this presupposition. The deployment posture proposed above, mandatory human review on every classification, sidesteps the calibration problem by removing the system's discretion to release outputs without review, at the cost of reducing the system's labour-saving value. The architectural redesign of the trigger, with which the system could plausibly operate with the automatic-review-by-exception posture originally intended, is identified as a precondition for full Article 14 compliance and is developed as future work in Section 8.1.

\hypertarget{comparison-with-state-of-the-art}{%
\section{Comparison with State-of-the-Art}\label{comparison-with-state-of-the-art}}

A direct quantitative comparison between the IDRISK2 system and the prior state of the art in automated FoodEx2 classification is constrained because the most advanced existing system, FEAST \parencite{lorenzo_molfetta_47d7c4cd}, operates on a fundamentally different input modality. FEAST consumes pre-structured textual food descriptions and produces multi-hierarchical FoodEx2 classifications through a retriever-reranker architecture; it does not ingest images. The 34-image dataset used in the present evaluation cannot be passed to FEAST without an upstream description-generation step, which is precisely the role that the OCR-and-VLM stages of the IDRISK2 system perform. The two systems address different points on the same end-to-end pipeline: FEAST is the textual-description-to-FoodEx2 specialist that the IDRISK2 system's RAG layer is broadly an instance of, while the IDRISK2 system as a whole is the image-to-FoodEx2 generalist that includes a FEAST-like component as its terminal stage.

The implication of this architectural separation is that the IDRISK2 system addresses the harder problem and accepts a correspondingly lower accuracy as the price of taking real-world inputs. The 29.4\% exact-match accuracy reported in Section 5.6 reflects the compounded error of the upstream VLM (which introduces visual misidentification, hallucinated codes, and contaminated queries) before any taxonomy-classification step has begun. The substantially higher accuracy reported by FEAST on structured-text inputs \parencite{lorenzo_molfetta_47d7c4cd} is not directly comparable because those figures measure the accuracy of the taxonomy-classification step in isolation, downstream of an assumed-correct textual description, whereas the IDRISK2 figure measures the end-to-end accuracy from raw image input. A faithful comparison would require either deploying FEAST on the VLM-generated descriptions produced by the IDRISK2 pipeline (which would introduce a stack-mismatch evaluation) or constructing a parallel dataset of textual descriptions that bypasses the visual stage entirely (which would not exercise the IDRISK2 system's core contribution).

The IDRISK2 system's specific contribution relative to the state of the art is therefore the demonstration of the end-to-end image-to-FoodEx2 pipeline on real-world dish photographs, complete with the closed-deployment, GDPR-compliant, EU AI Act-aware infrastructure that the literature reviewed in Section 3.5 identified as missing from existing approaches. The accuracy of the system is lower than what a structured-text classifier could achieve on cleaner inputs, but the system addresses an operational scenario, point-of-sampling classification by inspectors using smartphone photographs, that no published prior system has covered. The four gaps identified in Section 3.5 (end-to-end image-to-FoodEx2 pipeline, VLM evaluation on FoodEx2 classification, regulatory compliance integration, and closed self-hosted deployment) are addressed by the IDRISK2 system, even where the accuracy of the present implementation falls short of operational sufficiency.

\hypertarget{limitations-of-the-current-study}{%
\section{Limitations of the Current Study}\label{limitations-of-the-current-study}}

The evaluation reported in Chapter 5 is subject to five limitations that must be carried forward into the interpretation of the findings and the planning of future work.

First, the dataset size of 34 images is small for an empirical evaluation of an AI system. The reported accuracy and calibration figures have wide statistical uncertainty: an exact-match accuracy of 10/34 has a 95\% confidence interval of approximately 16\% to 47\%, and the 4 label-leak cases observed in the evaluation could plausibly correspond to a population rate anywhere between 3\% and 27\%. The qualitative findings, that the VLM hallucinates structurally, that the RAG layer is the source of all useful accuracy, that the calibration signal is uninformative, are robust to this uncertainty because they arise from patterns rather than precise percentages, but the specific numerical claims should be interpreted as point estimates with wide intervals. A larger evaluation set, ideally several hundred images, would be required to fix the metrics with the precision that would support deployment decisions.

Second, the ground truth was assigned by a single annotator. Inter-annotator agreement on FoodEx2 classification of culturally specific dishes is itself an open question, and the present evaluation cannot distinguish between cases where the system disagreed with the ground truth and cases where two annotators might have disagreed about the ground truth. The decision to record the system's family-level accuracy alongside its exact-match accuracy was motivated in part by this concern: the family-level metric is robust to disagreements between annotators about the specific sub-type within an agreed-upon family. Future work should include a multi-annotator validation on at least a subset of the evaluation set, with Cohen's Kappa or Fleiss's Kappa reported as a measure of the agreement ceiling against which the system's accuracy should be evaluated.

Third, the dataset is biased toward Portuguese and Brazilian dishes, which oversamples the cultural-coverage gap of the FoodEx2 taxonomy discussed in Section 6.3. A more representative dataset spanning the full geographical scope of EFSA's regulatory remit would be needed to estimate the system's accuracy in the operational European context. The bias is, however, deliberate for the purposes of the present evaluation, which aimed specifically to expose the system to inputs that would stress its coverage and grounding properties. The reported accuracy figures should therefore be read as a lower-bound estimate of what the system would achieve on a European-balanced dataset, not as a representative central estimate.

Fourth, the VLM evaluation reported in this dissertation covered all three interchangeable VLM backends supported by the system (Section 4.4.1): GPT-4o (OpenAI), Claude Sonnet 4.6 (Anthropic), and LLaVA-NeXT-Mistral-7B \parencite{liu_llava_next_2024}, all reported in Section 5.7.2. The included open-source evaluation establishes the boundary of the substitutability claim of Section 6.2 at the proprietary-to-7B-open-source transition and quantifies the accuracy collapse that results, but the open question of whether other contemporary open-source VLMs of comparable scale would reproduce the same failure profile, or whether the LLaVA-specific generic-fallback regime is idiosyncratic, remains formally open. A further open question is whether an instruction-tuned open-source VLM at a larger parameter count, or one fine-tuned on FoodEx2 documentation and labelled food imagery, could reduce the hallucination rate observed at 7B scale at the price of higher local compute cost. The cross-VLM evidence reported in Section 5.7.2 already supports the architectural substitutability claim within the proprietary class and refutes it at the proprietary-to-7B-open-source boundary; broader confirmation against additional open-source families and larger parameter counts is identified as follow-up work in Section 8.1.

Fifth, the efficiency criterion of the central research question (Section 1.3) is only partially addressed by the present evaluation. The system's per-image latency was measured (median 22.3 seconds, standard deviation 1.5 seconds) and the reduction of the inspector's effective search space from 31,303 taxonomy terms to a five-candidate reranked set was quantified, both of which characterise the system's intrinsic computational behaviour. What was not measured, and what the efficiency criterion strictly requires, is a comparative A/B evaluation against the current manual classification workflow with quantified end-to-end inspector throughput, time-to-classification, and cognitive load. The claim that the system improves the efficiency of FoodEx2 coding therefore rests on the latency and search-space-reduction indicators as preconditions for efficiency gains, rather than on a directly measured efficiency differential. The A/B field study required to quantify the efficiency differential is identified as a longitudinal-study priority in Section 8.5.
